\documentclass[11pt]{article}
\usepackage[utf8]{inputenc}
\usepackage{amsmath,amsthm,amssymb}
\usepackage{hyperref}
\usepackage{booktabs}
\usepackage{listings}
\usepackage[margin=1in]{geometry}

\newtheorem{theorem}{Theorem}
\newtheorem{corollary}[theorem]{Corollary}
\newtheorem{lemma}[theorem]{Lemma}
\newtheorem{definition}[theorem]{Definition}
\newtheorem{observation}[theorem]{Observation}

\title{Fibonacci Representation of Rational Numbers}
\author{Jan Popelka}
\date{December 2025}

\begin{document}

\maketitle

\begin{abstract}
We present a universal representation theorem: every rational number $p/q$ can be written as a sum of Fibonacci fractions $(\sum_i F_{a_i})/F_n$, where $n$ is the Fibonacci entry point of $q$ and $\{a_i\}$ is the Zeckendorf decomposition of $p \cdot (F_n/q)$. This extends classical Zeckendorf representation from integers to all rationals. As a corollary, we establish a polynomial bijection: every rational corresponds to a unique factored polynomial expression $P(x)/Q(x) \in \mathbb{Z}[x]/\mathbb{Z}[x]$ such that $P(\varphi)/Q(\varphi) = p/q$, where $\varphi$ is the golden ratio. We explain why this construction works using Galois theory of $\mathbb{Q}(\sqrt{5})$: the anti-symmetric form $\varphi^k - \psi^k$ produces $\sqrt{5}$ factors that cancel in the ratio.
\end{abstract}

\section{Introduction}

Zeckendorf's theorem \cite{zeckendorf1972} states that every positive integer has a unique representation as a sum of non-consecutive Fibonacci numbers. The Fibonacci entry point (or rank of apparition) $z(q)$ is the smallest positive integer $n$ such that $q$ divides $F_n$---a value guaranteed to exist by the Pisano period property.

Freitag and Filipponi \cite{freitag1990,filipponi1993} studied Zeckendorf representations of specific Fibonacci ratios $F_{kn}/F_n$. We generalize this to \emph{all} rational numbers by combining the entry point with Zeckendorf decomposition.

\textbf{Main contributions:}
\begin{enumerate}
    \item Universal representation theorem for $\mathbb{Q}$
    \item Polynomial bijection $\mathbb{Q} \leftrightarrow \mathbb{Z}[x]/\mathbb{Z}[x]$
    \item Extension to rational approximation of irrationals
    \item Efficient algorithmic implementation
\end{enumerate}

\section{Preliminaries}

\begin{definition}[Fibonacci Entry Point]
For $q \in \mathbb{Z}^+$, the entry point $z(q)$ is the smallest $n \geq 1$ such that $q \mid F_n$.
\end{definition}

\begin{theorem}[Pisano]
For every $q \in \mathbb{Z}^+$, the entry point $z(q)$ exists and divides the Pisano period $\pi(q)$.
\end{theorem}

\begin{definition}[Zeckendorf Representation]
For $n \in \mathbb{Z}^+$, the Zeckendorf representation is the unique set of indices $\{a_1 > a_2 > \cdots > a_k\}$ with $a_i - a_{i+1} \geq 2$ such that $n = \sum_i F_{a_i}$.
\end{definition}

\section{Main Results}

\subsection{Universal Representation Theorem}

\begin{theorem}[Fibonacci Fraction Representation]
\label{thm:main}
Every rational $p/q$ in lowest terms can be written as:
\[
\frac{p}{q} = \frac{\sum_{i} F_{a_i}}{F_n}
\]
where:
\begin{itemize}
    \item $n = z(q)$ is the Fibonacci entry point of $q$
    \item $\{a_i\} = \mathrm{Zeck}(p \cdot F_n/q)$ is the Zeckendorf decomposition of the scaled numerator
\end{itemize}
\end{theorem}

\begin{proof}
Since $z(q)$ is the entry point, we have $q \mid F_n$, so $F_n/q \in \mathbb{Z}$. Let $m = p \cdot (F_n/q)$. Then:
\[
\frac{p}{q} = \frac{p \cdot (F_n/q)}{F_n} = \frac{m}{F_n} = \frac{\sum_i F_{a_i}}{F_n}
\]
where the last equality uses Zeckendorf's theorem on $m$.
\end{proof}

\begin{corollary}[GCD Property]
For reduced $p/q$:
\[
\gcd\left(\sum_i F_{a_i}, F_n\right) = \frac{F_n}{q}
\]
\end{corollary}

\subsection{Polynomial Bijection}

\begin{theorem}[$\varphi$-Polynomial Representation]
\label{thm:poly}
Every rational $p/q$ corresponds to a unique factored polynomial expression:
\[
\frac{p}{q} = \frac{P(\varphi)}{Q(\varphi)}
\]
where $P(x), Q(x) \in \mathbb{Z}[x]$ are derived from the Fibonacci fraction representation via Binet's formula.
\end{theorem}

\textbf{Construction.} Using Binet's formula $F_n = (\varphi^n - \psi^n)/\sqrt{5}$ where $\psi = 1 - \varphi$:
\[
\frac{p}{q} = \frac{\sum_i (\varphi^{a_i} - \psi^{a_i})}{\varphi^n - \psi^n}
\]
The $\sqrt{5}$ factors cancel. Substituting $x$ for $\varphi$ and $(1-x)$ for $\psi$:
\[
\frac{P(x)}{Q(x)} = \frac{\sum_i (x^{a_i} - (1-x)^{a_i})}{x^n - (1-x)^n}
\]

\textbf{Example.} For $7/11$:
\begin{itemize}
    \item Entry point: $z(11) = 10$
    \item Scaled numerator: $7 \cdot (55/11) = 35$
    \item $\mathrm{Zeckendorf}(35) = \{9, 2\}$ since $35 = 34 + 1 = F_9 + F_2$
    \item Representation: $7/11 = (F_9 + F_2)/F_{10} = (34 + 1)/55$
\end{itemize}

Polynomial form:
\[
\frac{7}{11} = \frac{(x^9 - (1-x)^9) + (x^2 - (1-x)^2)}{x^{10} - (1-x)^{10}}
\]

\subsection{Special Point $x = 1/2$}

\begin{observation}
At $x = 1/2$, we have $x = 1-x$, so:
\begin{itemize}
    \item Every term $x^n - (1-x)^n = 0$
    \item Both $P(1/2) = 0$ and $Q(1/2) = 0$
    \item The ratio $P(x)/Q(x)$ has removable singularity at $x = 1/2$
\end{itemize}
\end{observation}

\textbf{Geometric note:} $\varphi - 1/2 = \sqrt{5}/2 \approx 1.118$.

\section{Algebraic Structure}

The representation theorem has a natural interpretation in terms of the algebraic number field $\mathbb{Q}(\sqrt{5})$.

\subsection{The Ring $\mathbb{Z}[\varphi]$}

The golden ratio $\varphi = (1+\sqrt{5})/2$ is an algebraic integer satisfying $\varphi^2 = \varphi + 1$. The ring of integers in $\mathbb{Q}(\sqrt{5})$ is:
\[
\mathbb{Z}[\varphi] = \{a + b\varphi : a, b \in \mathbb{Z}\}
\]

Every power of $\varphi$ reduces to this form via the recurrence $\varphi^n = F_n \cdot \varphi + F_{n-1}$, where $F_n$ is the $n$-th Fibonacci number.

\subsection{Galois Conjugation}

The field $\mathbb{Q}(\sqrt{5})$ has Galois group $\mathrm{Gal}(\mathbb{Q}(\sqrt{5})/\mathbb{Q}) = \{1, \sigma\}$ where the conjugation $\sigma$ acts by:
\[
\sigma(\varphi) = \psi = 1 - \varphi = \frac{1-\sqrt{5}}{2}
\]

The key Galois-theoretic operations are:
\begin{itemize}
    \item \textbf{Trace:} $\mathrm{Tr}(\varphi^n) = \varphi^n + \psi^n = L_n$ (Lucas numbers)
    \item \textbf{Norm:} $\mathrm{N}(\varphi^n) = \varphi^n \cdot \psi^n = (\varphi\psi)^n = (-1)^n$
\end{itemize}

\subsection{Why $\sqrt{5}$ Cancels}

The polynomial $f(x, k) = x^k - (1-x)^k$ satisfies the \emph{anti-symmetry property}:
\[
f(1-x, k) = -f(x, k)
\]

Under Galois conjugation $\sigma: \varphi \mapsto \psi = 1-\varphi$:
\[
f(\varphi, k) = \varphi^k - \psi^k = F_k \cdot \sqrt{5}
\]

This is precisely the Binet formula. In our rational representation:
\[
\frac{P(\varphi)}{Q(\varphi)} = \frac{\sum_i (\varphi^{a_i} - \psi^{a_i})}{\varphi^n - \psi^n} = \frac{\sum_i F_{a_i} \cdot \sqrt{5}}{F_n \cdot \sqrt{5}} = \frac{\sum_i F_{a_i}}{F_n} \in \mathbb{Q}
\]

The $\sqrt{5}$ factors cancel because both numerator and denominator use the \emph{anti-symmetric} combination $\varphi^k - \psi^k$, which always produces $F_k \cdot \sqrt{5}$.

\begin{theorem}[Galois Invariance]
For any $p/q \in \mathbb{Q}$, the polynomial representation satisfies:
\[
\frac{P(\varphi)}{Q(\varphi)} = \frac{P(\psi)}{Q(\psi)} = \frac{p}{q}
\]
That is, the value is fixed by Galois conjugation, confirming it lies in $\mathbb{Q}$.
\end{theorem}

\subsection{Fibonacci vs Lucas}

The Fibonacci and Lucas sequences arise from complementary Galois operations:

\begin{center}
\begin{tabular}{@{}lll@{}}
\toprule
Sequence & Formula & Galois Property \\
\midrule
$F_n$ & $(\varphi^n - \psi^n)/\sqrt{5}$ & Anti-symmetric (involves $\sqrt{5}$) \\
$L_n$ & $\varphi^n + \psi^n$ & Symmetric (trace, no $\sqrt{5}$) \\
\bottomrule
\end{tabular}
\end{center}

A ``Lucas fraction'' representation would use $g(x,k) = x^k + (1-x)^k$, which is symmetric: $g(1-x,k) = g(x,k)$. However, Lucas numbers lack a universal entry point property: not every integer divides some Lucas number.

\begin{observation}
The Fibonacci representation works precisely because:
\begin{enumerate}
    \item Every $q$ divides some $F_n$ (Pisano period property)
    \item The anti-symmetric form $\varphi^k - \psi^k$ produces $\sqrt{5}$, which cancels
    \item The result is Galois-invariant, hence rational
\end{enumerate}
\end{observation}

\section{Complexity Analysis}

\begin{definition}[Fibonacci Complexity]
The complexity $C(p/q)$ of a rational is the pair $(n, k)$ where $n = z(q)$ is the entry point and $k$ is the number of Zeckendorf terms.
\end{definition}

\begin{table}[h]
\centering
\begin{tabular}{@{}lcrl@{}}
\toprule
$p/q$ & Entry Point & Terms & Notes \\
\midrule
$1/2$ & 3 & 1 & Simplest non-trivial \\
$7/11$ & 10 & 2 & Moderate \\
$22/7$ & 8 & 3 & $\pi$ approximation \\
$355/113$ & 19 & 7 & Better $\pi$ approximation \\
$127/8191$ & 8190 & 2245 & Mersenne-related, very complex \\
\bottomrule
\end{tabular}
\caption{Fibonacci complexity of selected rationals}
\end{table}

\begin{observation}
Primes $q$ with large $z(q)$ (near $2q$) yield high complexity. Mersenne primes are particularly expensive.
\end{observation}

\section{Approximation of Irrationals}

\textbf{Algorithm} (FibonacciRationalize). Given $x \in \mathbb{R}$ and accuracy $\varepsilon$:
\begin{enumerate}
    \item For $n = 3, 4, 5, \ldots$:
    \item \quad $m = \mathrm{Round}(x \cdot F_n)$
    \item \quad If $|m/F_n - x| < \varepsilon$: return FibonacciFraction$(m/F_n)$
\end{enumerate}

\textbf{Example.} $\pi$ with accuracy $10^{-6}$:
\begin{itemize}
    \item Best rational: $355/113$ (entry point 19, 7 terms)
    \item Polynomial: degree 20 numerator, degree 18 denominator
\end{itemize}

\section{Related Work}

\begin{itemize}
    \item \textbf{Zeckendorf (1972)}: Original theorem for integers
    \item \textbf{Freitag \& Filipponi (1990, 1993)}: Zeckendorf representation of $F_{kn}/F_n$ ratios
    \item \textbf{Pisano periods}: Wall (1960), Wrench (1969)
    \item \textbf{Entry point properties}: Cubre \& Rouse (2012), Marques (2012)
    \item \textbf{Golden ratio base}: Bergman (1957)
\end{itemize}

Our contribution extends Freitag \& Filipponi from the specific subfamily $\{F_{kn}/F_n : k, n \in \mathbb{Z}^+\}$ to \emph{all} rationals $\mathbb{Q}$.

\section{Open Questions}

\begin{enumerate}
    \item \textbf{Polynomial structure}: What properties of $p/q$ are reflected in the factorization of $P(x)/Q(x)$?

    \item \textbf{Minimal representations}: Are there alternative Fibonacci fraction representations with fewer terms?

    \item \textbf{Complexity bounds}: Can we characterize rationals with bounded Fibonacci complexity?

    \item \textbf{Special point $x = 1/2$}: What is the limit $\lim_{x \to 1/2} P(x)/Q(x)$ and does it have meaning?

    \item \textbf{Lucas extension}: Does an analogous representation exist using Lucas numbers?

    \item \textbf{Connection to Egyptian fractions}: Both represent rationals as sums; what is their relationship?

    \item \textbf{Cyclotomic factors}: The sixth cyclotomic polynomial $(1-x+x^2)$ appears as a factor in $P(x)$ for certain rationals (e.g., $2/3$, $2/5$, $8/3$, $16/5$). At $x = \varphi$, this factor evaluates to $2$. What determines when cyclotomic factors appear, and how do they relate to Pisano periods?
\end{enumerate}

\section{Implementation}

A complete Wolfram Language implementation is available in the Orbit paclet:\footnote{Source code: \url{https://github.com/jan/orbit} (FibonacciFractions.wl)}

\begin{lstlisting}[language=Mathematica,basicstyle=\ttfamily\small]
<< Orbit`

(* Basic representation *)
FibonacciFraction[7/11]
(* {10, {9, 2}} *)

(* Polynomial form *)
FibonacciFraction[7/11, Method -> "Polynomial"]
(* Factored P(x)/Q(x) *)

(* Approximate irrationals *)
FibonacciRationalize[Pi, 10^-6]
(* {19, {21, 17, 14, 12, 10, 7, 4}} *)
\end{lstlisting}

Available methods: \texttt{Indices}, \texttt{Expression}, \texttt{Sum}, \texttt{Terms}, \texttt{Count}, \texttt{GoldenRatio}, \texttt{Polynomial}, \texttt{Matrix}.

\begin{thebibliography}{99}

\bibitem{zeckendorf1972}
Zeckendorf, E. (1972). Représentation des nombres naturels par une somme de nombres de Fibonacci ou de nombres de Lucas. \textit{Bulletin de la Société Royale des Sciences de Liège}, 41, 179--182.

\bibitem{freitag1990}
Freitag, H.T. (1990). On the Representation of $\{F_{kn}/F_n\}$, $\{F_{kn}/L_n\}$, $\{L_{kn}/L_n\}$, and $\{L_{kn}/F_n\}$ as Zeckendorf Sums. In \textit{Applications of Fibonacci Numbers}, Vol.~3, pp.~107--114. Kluwer.

\bibitem{filipponi1993}
Filipponi, P. \& Freitag, H.T. (1993). The Zeckendorf Representation of $\{F_{kn}/F_n\}$. In \textit{Applications of Fibonacci Numbers}, Vol.~5, pp.~217--219. Kluwer.

\bibitem{wall1960}
Wall, D.D. (1960). Fibonacci series modulo $m$. \textit{American Mathematical Monthly}, 67(6), 525--532.

\bibitem{cubre2012}
Cubre, P. \& Rouse, J. (2012). Divisibility properties of the Fibonacci entry point. \textit{Proceedings of the AMS}. arXiv:1212.6221.

\bibitem{bergman1957}
Bergman, G. (1957). A number system with an irrational base. \textit{Mathematics Magazine}, 31(2), 98--110.

\end{thebibliography}

\vspace{1em}
\noindent\textbf{Acknowledgments:} Computational exploration assisted by Claude (Anthropic).

\end{document}
